\section{Diskussion}
\label{chap:sicherheit}

Predictionsmärkte sind nutzerfreundlich, weil sie keine tiefen technischen
Vorkenntnisse erfordern. Damit bieten sie eine niedrige Einstiegsschwelle.
Gleichzeitig zeigen sie, dass Ergebnisfindung, Transparenz und
Informationsasymmetrien zentrale Fragen sind. Die DApp bildet diesen Spannungs-
bogen in vereinfachter Form ab.

Die Implementierung ist ein Demonstrationssystem und nicht produktionsreif.
Die Auflösung erfolgt admin-basiert. Eine zentrale Instanz bestimmt das Ergebnis
nach Fristende und erzeugt so einen Vertrauenspunkt. Im vorliegenden System
fehlen ein dezentrales Oracle- oder Streitbeilegungsmodell, sekundärer Handel
und weitergehende Absicherungen gegen missbräuchliche Markterstellung.

Auch Gas- und Benutzungsaspekte bleiben reduziert. Die Anwendung zeigt nur die
für den Kernfluss notwendigen Transaktionen und vermeidet zusätzliche
Komplexität. Das begrenzt Funktionsumfang und Realitätsnähe.

Ein dezentrales Oracle- oder Streitbeilegungsmodell reduziert die Abhängigkeit
von einer einzelnen Admin-Adresse. Nutzer können Ergebnisse vorschlagen. Die
Plattform vermittelt bei Konflikten oder Anfechtungen. Solche Konzepte finden
sich in produktiven Protokollen wie Polymarket.

Weitere Erweiterungen betreffen Marktvielfalt, Auswertungen zur Gas-Nutzung und
eine klarere Trennung zwischen UI-Vereinfachungen und on-chain-Allgemeingültigkeit.
Die aktuelle Lösung bildet einen didaktisch klaren Kern ab, der technisch
nachvollziehbar bleibt.

Ein dezentrales Oracle- oder Streitbeilegungsmodell reduziert die Abhängigkeit
von einer einzelnen Admin-Adresse. Nutzer können Ergebnisse vorschlagen. Die
Plattform vermittelt bei Konflikten oder Anfechtungen. Solche Konzepte finden
sich in produktiven Protokollen wie Polymarket.

Weitere Erweiterungen betreffen Marktvielfalt, Auswertungen zur Gas-Nutzung und
eine klarere Trennung zwischen UI-Vereinfachungen und on-chain-Allgemeingültigkeit.
Die aktuelle Lösung bildet einen didaktisch klaren Kern ab, der technisch
nachvollziehbar bleibt.